\documentclass[11pt,a4paper]{article}
\usepackage[T1]{fontenc}
\usepackage[utf8]{inputenc}
\usepackage{newtxtext,newtxmath}
\usepackage[margin=23mm,headheight=14pt,headsep=8mm,footskip=11mm]{geometry}
\usepackage{microtype}
\usepackage{amsmath}
\usepackage{booktabs,tabularx,array}
\usepackage{enumitem}
\usepackage{xcolor}
\usepackage{tikz}
\usetikzlibrary{arrows.meta,patterns}
\usepackage{caption}
\usepackage{titlesec}
\usepackage{fancyhdr}
\usepackage{xurl}

\usepackage[numbers,square,sort&compress]{natbib}
\bibliographystyle{unsrtnat}
\usepackage[
  colorlinks=true,linkcolor=accent,citecolor=accent,urlcolor=accent,
  pdftitle={Cherenkov diffraction radiation: analytical models, FDTD interfaces and an IPPL proof of concept},
  pdfsubject={Single-electron radiation at a vacuum-dielectric interface and numerical validation}
]{hyperref}

\definecolor{ink}{HTML}{173448}
\definecolor{accent}{HTML}{176B87}
\definecolor{pale}{HTML}{EAF2F5}
\definecolor{gridgray}{HTML}{CEDAE0}
\titleformat{\section}{\large\bfseries\color{ink}}{\thesection}{0.65em}{}
\titleformat{\subsection}{\normalsize\bfseries\color{ink}}{\thesubsection}{0.65em}{}
\titlespacing*{\section}{0pt}{16pt}{7pt}
\titlespacing*{\subsection}{0pt}{11pt}{5pt}
\setlength{\parindent}{0pt}
\setlength{\parskip}{5pt plus 1pt}
\setlength{\emergencystretch}{2em}
\setlist{nosep,leftmargin=1.5em,topsep=5pt}
\captionsetup{font=small,labelfont=bf,skip=7pt}
\pagestyle{fancy}
\fancyhf{}
\fancyhead[L]{\small\color{ink}Cherenkov diffraction radiation}
\fancyhead[R]{\small Literature and modelling note}
\fancyfoot[C]{\small\thepage}
\renewcommand{\headrulewidth}{0.3pt}
\newcolumntype{Y}{>{\raggedright\arraybackslash}X}
\newcommand{\dd}{\mathrm d}
\newcommand{\unit}[1]{\,\mathrm{#1}}
\newcommand{\code}[1]{\texttt{\small #1}}
\numberwithin{equation}{section}
\setlength{\bibsep}{3pt}
\renewcommand*{\bibfont}{\small}

\author{A Adelmann}
\title{Cherenkov diffraction radiation}
 
\begin{document}
\maketitle
\thispagestyle{plain}

This report illustrates a simple ChDR setup, the single-electron
literature review, the review of finite-difference time-domain (FDTD) material
interfaces, and the implications for the {\tt ChDR} mini-app based on IPPL FEL. The electron
remains in vacuum at a gap \(a\) from the radiator. Entering the longitudinal
region beside the material is distinct from crossing into the dielectric.

\begin{figure}[!ht]
\centering
\begin{tikzpicture}[
  x=1cm,
  y=1cm,
  >=Latex,
  font=\sffamily,
  beam/.style={very thick, -{Latex[length=3mm]}},
  radiation/.style={very thick, -{Latex[length=2.5mm]}, green!45!black},
  boundary/.style={very thick, black},
  gridline/.style={thin, gray!55},
  label/.style={font=\small},
  every node/.style={inner sep=2pt}
]

  % Vacuum background and dielectric radiator.
  \fill[blue!3] (-1.0,-2.8) rectangle (11.1,3.0);
  \fill[orange!18] (4.0,-2.45) rectangle (8.65,0.0);
  \draw[boundary] (4.0,0.0) -- (8.65,0.0) -- (8.65,-2.45)
                  -- (4.0,-2.45) -- cycle;

  % Computational mesh spanning both vacuum and dielectric.
  \foreach \xx in {-0.55,-0.10,0.35,0.80,1.25,1.70,2.15,2.60,3.05,3.50,
                    3.95,4.40,4.85,5.30,5.75,6.20,6.65,7.10,7.55,8.00,
                    8.45,8.90,9.35,9.80,10.25}
    \draw[gridline] (\xx,-2.45) -- (\xx,2.65);
  \foreach \yy in {-2.05,-1.65,-1.25,-0.85,-0.45,-0.05,0.35,0.75,
                    1.15,1.55,1.95,2.35}
    \draw[gridline] (-0.75,\yy) -- (10.65,\yy);
  \draw[gridline,dashed] (-0.75,-2.45) rectangle (10.65,2.65);

  % Beam travelling in physical +z (horizontal in TikZ), parallel to the surface.
  \draw[beam] (-0.75,0.72) -- (10.65,0.72);

  % A small bunch marker and negative charge sign.
  \draw[thick,fill=white] (1.05,0.72) circle (0.25);
  \node at (1.05,0.72) {$e^{-}$};

  % Beam--surface impact parameter.
  \draw[<->] (3.30,0.02) -- node[right,fill=white] {$a$} (3.30,0.70);
  \node[label,anchor=east,fill=white] at (6.4,0.4) {impact parameter};

  % Radiation generated in the dielectric at the Cherenkov angle.
  \draw[radiation] (4.18,0.0) -- (5.08,-0.957);
  \draw[radiation] (4.85,0.0) -- (5.75,-0.957);
  \draw[radiation] (5.52,0.0) -- (6.42,-0.957);
  \draw[radiation] (6.19,0.0) -- (7.09,-0.957);
  \draw[radiation] (6.86,0.0) -- (7.76,-0.957);
  \draw[radiation] (7.53,0.0) -- (8.43,-0.957);

  % Qualitative finite-edge radiation; not a Fresnel exit ray.
  \draw[radiation,dashed] (8.65,-2.45) -- (9.55,-1.25);
  \node[label,anchor=west,align=left,green!35!black] at (9.58,-1.25)
    {edge radiation\\(schematic)};

  % Cherenkov angle, measured from the beam direction (+z).
  \draw[thin] (4.18,0.0) ++(-46.75:0.55)
    arc[start angle=-46.75,end angle=0,radius=0.55];
  \node[label,anchor=west,fill=orange!18] at (4.30,-0.22)
    {$\theta_{\rm Ch}\simeq46.7^\circ$};

  % Material labels.
  \node[label,align=center] at (6.20,-1.55)
    {dielectric radiator\\
     $\epsilon_r=2.13$\;(illustrative)};
  \node[label,align=center] at (1.50,2.15)
    {vacuum\\$\epsilon_r=1$};

  % Coordinate axes.
  \draw[->,thick] (-0.55,-2.55) -- (-0.55,-1.65)
    node[left] {$x$};
  \draw[->,thick] (-0.55,-2.55) -- (0.35,-2.55)
    node[below] {$z$};

\end{tikzpicture}
\caption{Schematic \(x\)--\(z\) view of the setup draft.
The electron travels in vacuum along \(+z\), at gap \(a\) above the finite
radiator. The computational grid encompasses both materials. Solid green arrows indicate
internal radiation; the dashed edge-radiation arrow is qualitative, not a
prediction of the external observation angle.}
\label{fig:geometry}
\end{figure}




%\textbf{Main finding.} The infinite planar-interface problem supplies an
%analytical reference requiring one numerical integral per frequency.
%A finite radiator additionally requires edge and extraction effects.
%FDTD can model both, provided the constitutive response, moving-charge source
%and outer boundary are consistent. A Cartesian grid can represent a tilted
%face; the current FEL solver needs a material formulation before that geometry
%can be used quantitatively.

\section{Geometry and fixed beam parameters}
\label{sec:setup}

Let us fix the AWA like beam parameters as follows, assuming \(60\unit{MeV}\) means
kinetic energy and \(3\unit{mm}\) means the rms longitudinal length of a
Gaussian bunch. The magnitude of the bunch charge is \(5\unit{nC}\).
The beam travels along \(+z\), with \(x\) normal to the beam-facing surface
and \(y\) the other transverse coordinate. Thus \(\sigma_z=3\unit{mm}\);
\(\sigma_x\) and \(\sigma_y\) denote the transverse beam sizes.
The dielectric (fused silika) is initially an illustrative nonmagnetic, lossless medium
with relative permittivity \(\epsilon_r=2.13\). 

\begin{center}
\renewcommand{\arraystretch}{1.14}
\begin{tabularx}{\linewidth}{@{}Y l@{}}
\toprule
Quantity & Value\\
\midrule
Electron kinetic energy \(K\) & \(60\unit{MeV}\)\\
Total energy per electron \(E_{\mathrm{tot}}\) & \(60.511\unit{MeV}\)\\
Lorentz factor \(\gamma=1+K/(m_ec^2)\) & \(118.4171\)\\
Normalized speed \(\beta=\sqrt{1-\gamma^{-2}}\) & \(0.99996434\)\\
Longitudinal rms length \(\sigma_z\) & \(3\unit{mm}\)\\
Rms duration \(\sigma_t=\sigma_z/(\beta c)\) & \(10.0073\unit{ps}\)\\
Gaussian pulse FWHM & \(23.6\unit{ps}\)\\
Gaussian peak current magnitude \(I_{\mathrm{peak}}\) & \(199\unit{A}\)\\
Charge magnitude \(Q\), electron number \(N=Q/e\) &
\(5\unit{nC},\quad 3.12\times10^{10}\)\\
Bunch kinetic energy \(NK\) & \(0.30\unit{J}\)\\
Refractive index \(n=\sqrt{2.13}\) & \(1.45945\)\\
\bottomrule
\end{tabularx}
\end{center}

The duration, peak current and energy follow from
\begin{align}
 t_{\mathrm{FWHM}}&=2\sqrt{2\ln2}\,\sigma_t,&
 I_{\mathrm{peak}}&=\frac{Q}{\sqrt{2\pi}\,\sigma_t},\\
 E_{\mathrm{tot}}&=K+m_ec^2,&
 E_{\mathrm{bunch,kin}}&=NK.
\end{align}
%Here \(Q\) is the positive charge magnitude; the electron bunch carries
%charge \(-Q\). The approximation \(\sigma_t\simeq\sigma_z/c\) gives
%\(10.0\unit{ps}\), consistent with the setup draft.

\subsection{First simulation case}

A straight \(60\unit{MeV}\), \(5\unit{nC}\) bunch passing parallel to a flat
\(\epsilon_r=2.13\) radiator is a useful first proof of concept. Varying the
gap \(a\) tests coupling, spectrum, polarization and the internal Cherenkov
angle before adding the tilted prism and its optical extraction.
The gap, transverse beam sizes and radiator dimensions remain parameters.

A longitudinal span of \(6\)--\(10\,\sigma_z\), or approximately
\(18\)--\(30\unit{mm}\), is a starting allowance for the bunch in a
laboratory-frame domain. Add the radiator, upstream and downstream vacuum
margins, radiation travel distance and absorbing layers. This bunch-span
estimate alone does not determine the full simulation volume or duration.

\subsection{Simplified experimental considerations}

The \(3\unit{mm}\) bunch length gives a longitudinal squared-form-factor
roll-off near \(16\unit{GHz}\), corresponding to centimetre wavelengths.
This motivates examining a microwave band, but the detected maximum also
depends on the radiator response and collection geometry. A flat slab
validates generation inside the dielectric; the prism and optical system
are needed to predict the collected signal. Section~\ref{sec:analytical}
separates the single-electron response, bunch coherence and finite-body effects.

The prism aperture models discussed below are most useful when face dimensions
greatly exceed wavelength. A centimetre-scale radiator is comparable to the
approximately \(19\unit{mm}\) vacuum wavelength at this bunch form-factor
scale, so that assumption must be checked. The appropriate material input is
\(\epsilon_r(\omega)\) over the chosen band; an optical dielectric constant
cannot be presumed valid at microwave frequencies.

\section{Analytical and reduced single-electron models}
\label{sec:analytical}

\subsection{Literature closest to the stated geometry}

Tyukhtin, Galyamin and Vorobev~\cite{tyukhtin2021invertedprism}
give the exact dielectric-half-space field as a single transverse-wave-number
integral in Section~III, Eqs.~(5)--(8). This is the most direct reference
for the planar benchmark. The later finite-prism construction in that paper
uses an aperture approximation, and should not be treated as an exact
solution for any finite body.

Shevelev and Konkov~\cite{shevelev2014dielectrictarget} calculate polarization
radiation from a charge outside a prismatic target, including target
orientation and extraction. Their target is infinite in one transverse
direction. The Fresnel extraction assumes an exit face much larger than
the wavelength and omits multiple internal rereflections; see p.~510.
Their gap is denoted by \(b\), whereas their \(a\) is a target dimension.

Curcio et al.~\cite{curcio2020bunchlength} connect ChDR theory, simulations
and bunch-length measurements. Section~III.A compares MAGIC simulations of
an axisymmetric hollow cone with an analytical radiator transfer function.
The modulated beam and conical geometry support the approach, but are
different from a planar single-electron benchmark.

\subsection{Meaning of a one-dimensional computation}

For the analytical reference, replace the finite radiator in
Fig.~\ref{fig:geometry} by a dielectric half-space with an infinitely extended
beam-facing plane at \(x=0\). Choose vacuum at \(x>0\) and dielectric at
\(x<0\). Cartesian vectors below are ordered as \((x,y,z)\). For a prescribed
electron trajectory at gap \(a>0\), moving along \(+z\),
\begin{align}
 \mathbf r_e(t)&=(a,0,vt), \label{eq:trajectory}\\
 \rho(\mathbf r,t)&=-e\,\delta(x-a)\delta(y)\delta(z-vt),&
 \mathbf J&=v\rho\,\hat{\mathbf z}. \label{eq:source}
\end{align}
Here \(v=\beta c>0\), so \(J_x=J_y=0\) and \(J_z=v\rho\).
With time convention \(e^{-i\omega t}\), the source fixes
\begin{equation}
 k_z=\frac{\omega}{v}.
 \label{eq:phase}
\end{equation}
Fourier transformation in \(y\) is useful because the ideal half-space is
uniform along \(y\): spatial modes \(e^{ik_y y}\) can be solved independently.
Motion along \(z\) already fixes \(k_z=\omega/v\). Thus, for each pair
\((\omega,k_y)\), only the dependence on \(x\) remains to be solved as a
one-dimensional boundary-value problem. For \(\mu_r=1\), the vacuum decay
constant and dielectric normal wave number obey
\begin{align}
 \kappa_0&=\sqrt{k_y^2+\frac{\omega^2}{\gamma^2v^2}}, \label{eq:kappa}\\
 k_{x,d}^{\,2}&=\epsilon_r(\omega)\frac{\omega^2}{c^2}
                  -\frac{\omega^2}{v^2}-k_y^2. \label{eq:kxd}
\end{align}
The subscript \(d\) in \(k_{x,d}\) means \emph{dielectric}: it labels the
wave-number component normal to the interface inside the dielectric,
not a distance or derivative. Equation~\eqref{eq:kxd} follows from the
dielectric dispersion relation
\[
 k_{x,d}^{\,2}+k_y^2+k_z^2
 =\epsilon_r(\omega)\frac{\omega^2}{c^2}
\]
after substituting \(k_z=\omega/v\).
The vacuum field is evanescent normal to the interface. Choose the
dielectric branch to carry energy away from the interface, or to decay
into passive lossy material.

Each spectral component is obtained by matching the two polarizations
at the interface. The electron is localized at \(y=0\), so its field
contains all \(k_y\) components even though the material is uniform in
\(y\). Integrating these components over \(k_y\) reconstructs the
three-dimensional point-electron field at a fixed frequency. Reconstructing
a time waveform also requires an inverse frequency transform. Both
propagating and evanescent dielectric components belong in the full integral.

This is a \emph{spectral reduction of a three-dimensional point-source
problem}. Ordinary 1D FDTD cannot resolve both the gap and the propagation
direction. Ordinary Cartesian 2D FDTD describes a source uniform in the
omitted coordinate, hence a line charge with per-unit-length normalization.
The explicit amplitudes and Fourier normalization in
Ref.~\cite{tyukhtin2021invertedprism} use Gaussian electromagnetic units;
convert source and field normalization consistently for an absolute
comparison in SI or IPPL units.

\subsection{Threshold, internal angle and gap dependence}

For a real, positive, nondispersive permittivity, a propagating dielectric
component requires
\begin{equation}
 k_y^2<\frac{\omega^2}{c^2}
       \left(\epsilon_r-\beta^{-2}\right).
\end{equation}
Consequently the threshold and internal phase-matching angle are
\begin{equation}
 n\beta>1,\qquad
 \cos\theta_{\mathrm{Ch}}=\frac{1}{n\beta},\qquad
 \theta_{\mathrm{Ch}}=46.7476^\circ
 \quad(\epsilon_r=2.13).
 \label{eq:angle}
\end{equation}
This angle is measured inside the homogeneous dielectric relative to
the beam direction \(+z\). It does not specify the camera angle after extraction.

At a fixed \(k_y\), moving the trajectory away from the interface supplies
the amplitude factor \(e^{-\kappa_0a}\), and the corresponding modal energy
factor \(e^{-2\kappa_0a}\). An integrated observable combines many \(k_y\)
values and need not follow a single exponential in \(a\). At \(k_y=0\),
\begin{equation}
 \kappa_0^{-1}=\frac{\gamma\beta\lambda_0}{2\pi}.
 \label{eq:coupling}
\end{equation}
For example, this field coupling length is approximately
\(9.4\unit{\mu m}\) at a vacuum wavelength of \(500\unit{nm}\).
It is neither a hard cutoff nor a required box dimension.

\subsection{Extension to the specified bunch}

For prescribed trajectories in a linear material, electron responses can
be superposed. If they differ only by independent arrival times and have
the same transverse coupling, expanding the squared sum of the amplitudes gives
\begin{align}
 \left\langle\frac{\dd W_N}{\dd\omega}\right\rangle
   &=\left[N+N(N-1)|F(\omega)|^2\right]
      \frac{\dd W_1}{\dd\omega}, \label{eq:bunch}\\
 |F(\omega)|^2&=\exp[-(\omega\sigma_t)^2]
              =\exp[-(\omega\sigma_z/v)^2],&
 N&=\frac{Q}{e}. \label{eq:formfactor}
\end{align}
A transverse distribution, including a distribution of gaps, requires
an amplitude average over those trajectories.

The Gaussian squared form factor reaches \(e^{-1}\) at
\begin{equation}
 f_*=\frac{1}{2\pi\sigma_t}=\frac{v}{2\pi\sigma_z}=15.904\unit{GHz},
 \qquad \lambda_{0,*}\simeq18.85\unit{mm}.
\end{equation}
This is a form-factor roll-off, \emph{not automatically a detected spectral
maximum}. The single-electron spectrum, finite-radiator response and
collection optics also contribute. A direct representation of individual
electrons does not remove the grid's spatial and temporal bandwidth limits.
For a fixed bunch shape and geometry, the coherent contribution scales
approximately as \(Q^2\). At visible wavelengths the Gaussian longitudinal
coherent term is strongly suppressed for a \(3\unit{mm}\) bunch. A visible-camera
prediction therefore needs the incoherent single-electron contribution,
or optical-scale bunching if coherent emission is intended.

\subsection{Finite radiator and the approximately 47-degree geometry}

Equation~\eqref{eq:phase} implies \(k_z>\omega/c\): the tangential wave
number exceeds the vacuum propagating-wave limit. Therefore a mode
coupled through the infinite planar interface cannot escape as propagating
radiation through that same interface. An inclined extraction face or a
finite edge changes this constraint. This is a phase-matching consequence
of the idealized model.

A rectangular radiator can demonstrate excitation inside the dielectric,
while producing a different external angular distribution from a prism.
Distinguish the orientation of the beam-facing surface, the extraction-face
normal and the internal angle in Eq.~\eqref{eq:angle}.

Entering and leaving the region beside a finite body breaks translation
symmetry along \(z\). Edge transients, interference and internal reflections
then require a finite-geometry treatment. If the material is still invariant
in \(y\), a Fourier transform in that coordinate gives a family of
two-dimensional problems, often called 2.5D. Finite width generally requires
a three-dimensional calculation.

\section{FDTD and the vacuum--dielectric boundary}
\label{sec:fdtd}

\subsection{Field equations and the material interface}

The Meep paper by Oskooi et al.~\cite{oskooi2010meep} describes a practical
FDTD implementation and includes moving-source Cherenkov radiation in
Fig.~5. Its homogeneous-medium, two-dimensional demonstration is a useful
starting point. The maintained source example is also available~\cite{meepMovingSource}.

A suitable material formulation, written here in SI units, is
\begin{align}
 \partial_t\mathbf D&=\nabla\times\mathbf H-\mathbf J_{\mathrm{free}},
 &\partial_t\mathbf B&=-\nabla\times\mathbf E, \label{eq:maxwell}\\
 \mathbf B&=\mu_0\mathbf H,
 &\mathbf D&=\epsilon_0\epsilon_r(\mathbf r)\mathbf E. \label{eq:constitutive}
\end{align}
For the initial lossless, nondispersive proof of concept,
\(\epsilon_r=1\) in vacuum and \(2.13\) in the radiator.
The electron's field polarizes the material; a separate prescribed
Cherenkov-radiation source is unnecessary.

With no free surface-charge sheet or free surface current, the physical
interface conditions are
\begin{equation}
 [\mathbf E_t]=[\mathbf H_t]=0,\qquad [D_n]=[B_n]=0.
 \label{eq:interface}
\end{equation}
Here square brackets denote the jump across the interface. The normal
components at the beam-facing plane \(x=0\) are \(D_x,B_x\), and the
tangential components are \(E_y,E_z,H_y,H_z\). The normal
electric field generally jumps; forcing every component of \(\mathbf E\)
to be continuous gives an incorrect boundary. Bound polarization charge
is represented by the constitutive relation.

In a Yee discretization, material coefficients must correspond to the
electric-component locations and their averaging volumes. Beginning with a
grid-aligned plane isolates constitutive and source errors from errors
in representing a sloping face.

\subsection{Cut cells and a tilted extraction face}

Farjadpour et al.~\cite{farjadpour2006subpixel} derive subpixel averaging
that removes the leading interface error for suitable geometries.
Kottke et al.~\cite{kottke2008interfaces}, Section~VI, explain the
discontinuity-aware treatment and its anisotropic extension.

For a cell cut by a locally planar interface between two isotropic,
nondispersive materials, let \(f\) be the dielectric fraction.
Using relative permittivities, the tangential and normal averages are
\begin{align}
 \epsilon_{\parallel}&=f\epsilon_d+(1-f)\epsilon_v,\\
 \epsilon_{\perp}&=\left(
    \frac{f}{\epsilon_d}+\frac{1-f}{\epsilon_v}
    \right)^{-1}.
\end{align}
For unit interface normal \(\hat{\mathbf n}\), these form the tensor
\begin{equation}
 \boldsymbol{\epsilon}_{\mathrm{eff}}
 =\epsilon_{\parallel}
   (\mathbf I-\hat{\mathbf n}\hat{\mathbf n}^{T})
 +\epsilon_{\perp}\hat{\mathbf n}\hat{\mathbf n}^{T}.
 \label{eq:tensor}
\end{equation}
At the beam-facing plane, \(\hat{\mathbf n}=\hat{\mathbf x}\), so in
\((x,y,z)\) order this is
\(\operatorname{diag}(\epsilon_{\perp},\epsilon_{\parallel},\epsilon_{\parallel})\).
Off-diagonal entries require consistent interpolation on a staggered
grid. This describes a sharp interface within a numerical cell, rather
than a deliberately thick physical transition layer. Corners and edges
can limit convergence below second order~\cite{kottke2008interfaces}.

A \(47^\circ\) face is therefore possible on a Cartesian grid.
The issue is accurate representation and convergence of the interface.
Arbitrary scalar averaging of the permittivity does not implement
Eq.~\eqref{eq:tensor}.

\subsection{Frequency-dependent response}

For dispersive material, introduce a causal polarization model, for example
auxiliary equations for Lorentz or Debye responses, with
\begin{equation}
 \mathbf D=\epsilon_0\epsilon_\infty\mathbf E+\mathbf P.
\end{equation}
The nondispersive averaging above is not a complete dispersive algorithm.
Meep smooths the instantaneous permittivity and documents limitations
for dispersive contributions~\cite{meepSubpixel}. Material loss, bandwidth
and any auxiliary-equation stability conditions need explicit choices.

\subsection{Moving charge, initialization and conservation}
\label{sec:source-consistency}

Esirkepov's deposition construction~\cite{esirkepov2001charge}
provides a reference for satisfying discrete continuity. Its discrete
operators must match those of the selected Maxwell solver:
\begin{equation}
 \frac{\rho^{n+1}-\rho^n}{\Delta t}
 +\nabla_h\cdot\mathbf J^{n+1/2}=0.
 \label{eq:continuity}
\end{equation}
Combined with consistent initial fields, this preserves the free-charge
Gauss constraint \(\nabla_h\cdot\mathbf D=\rho\).
Starting an already moving electron with zero fields, or abruptly switching
its current on and off, creates transients. Initialize consistently or
show that those transients do not enter the benchmark measurement.

The numerical particle shape must fit within the vacuum gap and resolve
the measured frequencies. Otherwise deposition overlaps the material or
filters the spectrum. These requirements still apply when every simulated
particle represents one electron.

\subsection{Outer absorption and numerical radiation}

A perfectly matched layer (PML) truncates the computational domain;
it is separate from the physical vacuum--dielectric interface.
Johnson's notes~\cite{johnson2021pml} explain the construction and its
limitations. Place it outside the observation region, and verify convergence
with its distance and thickness. If dielectric reaches the PML, extend
the material consistently. Inhomogeneity along the stretching direction,
including some oblique structures, can invalidate ordinary
matching~\cite{meepPML}.

The setup draft additionally cites Kraus, Adelmann and
Arbenz~\cite{kraus2011adipml}, who formulate PML for a divergence-preserving
alternating-direction implicit (ADI) Maxwell scheme. This is relevant when
evaluating an ADI route; its update equations must not be assumed to apply
unchanged to the present explicit potential solver.

Xu et al.~\cite{xu2020numericalfields} analyse spurious fields and
Cherenkov-like wakes around a relativistic particle, including in vacuum.
At the present \(\beta=0.99996434\), a vacuum control with the same source,
mesh and time step is essential. A visible cone alone is not evidence
of physical ChDR. Numerical radiation from a prescribed charge should also
be distinguished from collective numerical Cherenkov instability.

Measure radiation separately from the electron's bound near field:
compare complex field spectra at specified locations, or outgoing flux
away from the trajectory. Align \(\mathbf E\) and \(\mathbf H\) in space
and time before computing Poynting flux~\cite{oskooi2010meep}.
Subtracting two near-field powers is not generally equivalent to computing
the power of a difference field.

\section{Implications for the new IPPL {\tt ChDR} mini-app}
\label{sec:ippl}

\subsection{Development scope for the MSc thesis of D.~Rosser}

The setup draft identifies dielectric interfaces, the reference frame and
absorbing boundaries as development topics for this thesis. The following
items connect that scope to the current implementation.

\textbf{Field formulation.}
\path{demos/fel/datatypes.h}, around line~50, selects
\code{NonStandardFDTDSolver} with absorbing Mur boundaries.
\path{src/MaxwellSolvers/NonStandardFDTDSolver.hpp}, around lines~28
and~114, advances four-potential fields using a vacuum stencil and sets
\(\Delta t=h_z\) in normalized units. A spatial dielectric is more than a
geometry mask or a replacement of \(c\) by \(c/n\): constitutive response,
interface coupling and stability must be derived consistently.

\textbf{Existing particle and source machinery.}
The FEL mini-app already provides relativistic Boris particle motion,
field interpolation, particle initialization and current deposition.
In \path{demos/fel/FreeElectronLaserManager.h}, \code{depositCurrent}
calls \code{assemble\_current\_collocated}; CIC charge-density deposition
is enabled by \code{space\_charge}. These are existing capabilities to
reuse for ChDR. The material extension does not by itself require a new
particle pusher or deposition scheme.

\textbf{Initialization and conservation checks.}
Section~\ref{sec:source-consistency} is therefore supported by existing
particle/source machinery, but its full consistency requirements have
not yet been established for the ChDR configuration.
In \path{src/MaxwellSolvers/NonStandardFDTDSolver.hpp},
\code{initialize} sets all three potential time levels to zero.
The manager's initialization sequence does not supply the pre-existing
self-field of an already moving electron. Consistent field initialization,
or exclusion of startup transients from the measurement, remains necessary.

The current implementation uses trajectory-segmented CIC deposition on
a collocated grid. The tests in
\path{unit_tests/Interpolation/TestCurrentDeposition.cpp} check deposited
values, but do not establish the discrete continuity relation in
Eq.~\eqref{eq:continuity} or preservation of the Gauss constraint.
Exact discrete charge conservation should therefore be verified using
the selected field operators; it is not established by the presence of
the deposition routine alone. Repeat that compatibility check if the
material solver or field staggering changes.

\textbf{Moving electron source:}
Oskooi and Johnson's chapter in the volume edited by Taflove, Oskooi and
Johnson~\cite{taflove2013advances,oskooi2013sources},
Section~4.7, pp.~89--91, explicitly demonstrates moving-source Cherenkov
radiation. Figure~4.10 compares continuous interpolation of the source
with rounding its position to grid points; the latter produces artificial
high-frequency radiation. Interpolation reduces these artifacts, while
numerical dispersion remains. This homogeneous-medium, two-dimensional
example supports source validation but does not replace the
vacuum-gap benchmark for our geometry. The chapter is authored by
Oskooi and Johnson; Taflove is a co-editor of the book.
Taflove and Hagness~\cite{taflove2005computational} provide the broader
FDTD background on material updates, numerical dispersion, stability
and absorbing boundaries.

\textbf{Reference frame.}
The same manager, around line~47, ties a Lorentz boost to bunch and
undulator parameters. A laboratory-frame ChDR mode keeps the radiator
stationary. A dielectric that is stationary and isotropic in the laboratory
cannot be represented as the same stationary scalar material after a
nonzero boost. Removing the undulator field alone does not establish
the correct material-frame model.

\textbf{Outer boundaries.}
The second-order Mur treatment can be used for preliminary checks if
reflections are shown to be below the measurement tolerance. Quantitative
power and oblique radiation motivate a validated PML implementation.
The absorber must remain separate from the physical dielectric interface.

Two possible formulations merit evaluation:
\begin{enumerate}
 \item A direct \(\mathbf D,\mathbf B\) Maxwell solver with a material
       constitutive response and consistent particle coupling.
 \item A potential solver coupled self-consistently to polarization
       charge and current,
       \begin{equation}
       \rho_{\mathrm{bound}}=-\nabla\cdot\mathbf P,\qquad
       \mathbf J_{\mathrm{bound}}=\partial_t\mathbf P.
       \end{equation}
       This route needs a stable, gauge-consistent coupling.
\end{enumerate}
The literature review does not establish which requires the smaller
implementation change.

\subsection{Job-file units and coordinate mapping}

The parser in \path{demos/fel/Config.h} converts the numerical
\code{charge} and \code{mass} entries from elementary-charge and
electron-mass units into solver units. They are not direct SI coulombs
and kilograms. The manager subsequently divides the configured totals
by the actual particle count.

The charge conversion factor in \path{demos/fel/units.h} is positive and
the input sign is passed through. Thus, for an electron bunch of charge
\(-5\unit{nC}\), the setup draft's charge magnitude becomes a negative
input. A sketch of the relevant entries is
\begin{verbatim}
"charge": -3.1207545e10,
"mass":    3.1207545e10,
"gamma":   118.417071,
"sigma-position": [sigma_x_code_um, sigma_y_code_um, 3000.0]
\end{verbatim}
The two transverse sizes are placeholders. The value \(3000.0\) assumes
the job file selects micrometres as its length scale. In the current code
the third coordinate, \(z_{\mathrm{code}}\), is longitudinal and corresponds
to \(z\) in Fig.~\ref{fig:geometry}; \(x_{\mathrm{code}}\) and
\(y_{\mathrm{code}}\) are transverse. Thus the third entry is
\(\sigma_z=3000\unit{\mu m}\), while the first two are \(\sigma_x,\sigma_y\).
This coordinate convention does not remove the code's boost.
These entries therefore specify units
and beam scales, not a complete runnable laboratory-frame ChDR job.

\subsection{Validation sequence}

\begin{enumerate}[itemsep=5pt]
 \item \textbf{Passive interface.} Launch plane waves at a flat boundary.
       Compare Fresnel reflection and transmission for both polarizations
       and check lossless energy balance.
 \item \textbf{Vacuum electron.} Use a prescribed \(60\unit{MeV}\)
       trajectory and consistent initial fields. Measure unwanted radiation,
       continuity residuals and the Gauss residual.
 \item \textbf{Planar ChDR.} Compare complex spectra, polarization, gap
       dependence and propagation angle with the independent half-space
       calculation. Check convergence in mesh spacing, particle shape,
       box dimensions and outer absorption.
 \item \textbf{Finite radiator.} Add leading and trailing edges, followed
       by the tilted extraction face. Compare angular spectra and refine
       the interface representation. Apply finite-target analytical
       approximations only within their stated regime.
 \item \textbf{Bunch response.} Introduce the longitudinal and transverse
       distributions in the selected band after the single-electron
       response is established.
\end{enumerate}

For an infinite track, compare a spectral energy quantity per unit path
length rather than a finite total radiated energy. For a finite radiator,
define whether the observable is a local field, electron energy loss,
outgoing spectral energy or energy collected by specified optics.

For a conventional explicit Yee solver, the vacuum region imposes
\begin{equation}
 c\Delta t\sqrt{h_x^{-2}+h_y^{-2}+h_z^{-2}}\leq1.
\end{equation}
Auxiliary material equations and other stencils can impose additional
conditions. This is not the time-step rule of the existing nonstandard
potential solver. Resolve the shortest material wavelength and relevant
gap/source scales; a prescribed cell count alone does not certify accuracy.

\subsection{Inputs still needed and status}

The remaining physical choices are gap \(a\), radiator length, thickness
and width, face orientations, material data \(\epsilon_r(\omega)\), the
frequency band, transverse beam distribution and the detector observable.
The approximately \(16\unit{GHz}\) bunch form-factor scale is useful for
planning, but does not determine the detected peak or eliminate the
need for the radiator transfer function.

This is a focused literature synthesis and implementation assessment.
The beam quantities are direct calculations from the stated assumptions.
The listed simulations and convergence checks are recommendations;
no field simulation or solver modification was performed for this report.

\clearpage
\renewcommand{\bibsection}{\section{Literature and software references}}
\bibliography{chdr_references}

\end{document}
